\documentclass[../Main.tex]{subfiles}

\begin{document}
\section{Superposition}
\begin{definition}{Boundary value problem}
    A \underline{boundary value problem} is a problem of the form:
    \begin{equation*}
        P\phi(\vec{x}) = 0
    \end{equation*}
    defined on a bounded, open subset $\Omega$ of $\R^n$. $\phi$ must also satisfy appropriate boundary conditions on $\partial \Omega$.
\end{definition}
\begin{definition}{Initial value problem}
    An \underline{initial value problem} is a problem of the form:
    \begin{equation*}
        P\phi(\vec{x}, t) = 0
    \end{equation*}
    defined on a bounded, open subset $\Omega$ of $\R^n$ and on $t > 0$. $\phi$ must satisfy appropriate boundary conditions on $\partial \Omega$ for all time, and appropriate initial conditions at $t = 0$ on $\Omega$.
\end{definition}
We can split IBVPs into a problem with initial condition $0$ (and with the given boundary conditions), and a problem with boundary conditions $0$ (and with the given initial conditions). The summing the solutions gives the required solution.

We can also split the problem into sub-problems with the boundary conditions only non-zero on one region of the boundary domain, such as only one side of a square.
\section{Laplace's Equation}
Laplace's equation is a boundary value problem given by:
\begin{equation*}
    \nabla^2 \phi = 0
\end{equation*}
where $\nabla^2 = \nabla \cdot \nabla$.

For Laplace's Equation on a square, disc or ball, we can transform to the appropriate coordinate system (Cartesian, cylindrical and spherical, respectively) and use separation of variables. Our general process is:
\begin{enumerate}
    \item Try a separable solution and solve for separate problems involving only one boundary condition
    \item Apply homogeneous boundary conditions to the individual eigenfunctions
    \item Multiply the separable parts back together, and sum over all possible eigenvalues (normally indexed by $\N$ or $\N_0$)
    \item Apply any non-homogeneous boundary conditions using the relevant inner product for the operator.
\end{enumerate}

We often find that $\lambda > 0$ is a requirement to avoid trivial solutions. This cannot be simply assumed, and often involves finding solutions for negative $\lambda$ and showing this cannot be non-trivial with the given boundary conditions.

For Laplace's Equation on a rectangle or disc, we expect solutions involving trigonometric functions (or hyperbolic trigonometric functions). For Laplace's Equation on a sphere, we must transform to $x = \cos(\theta)$ and find Legendre's Equation in $x$.
\section{Wave Equation}
\subsection{Solution}
By considering the tension in a taut string, we derive the Wave Equation:
\begin{equation*}
    \frac{\partial^{2}\phi}{\partial t^{2}} = c^2 \nabla^2 \phi
\end{equation*}
which is an IBVP.

We again consider separation of variables to solve this on a line or a disc. The solution $y(x, t)$ for a line with fixed ends at $0, L$ is:
\begin{equation*}
    y(x, y) = \sum_{n=0}^{\infty} \sin\left(\frac{n\pi x}{L}\right) \left[A_n \cos\left(\frac{n\pi ct}{L}\right) + B_n \sin\left(\frac{n\pi c t}{L}\right)\right]
\end{equation*}
where the $A_n$ and $B_n$ are determined by the initial conditions (on $y$ and $\dot{y}$).

The separation of variables on a disc results in Bessel's Equation with $m = 0$ for $R$, and a sinusoidal solution for $T$.
\subsection{Energy Conservation and Uniqueness}
For the first example with a string clamped at $0$ and $L$, define the energy:
\begin{equation*}
    E(t) = \int_{0}^{L} \left[\frac12 \mu y_t^2 + \frac12 \tau y_x^2\right] dx
\end{equation*}
where $\mu$ is the mass per unit length of the string, $\tau$ is the tension. We can set $\mu = 1$ so that $\tau = c^2$. Then differentiating gives that the energy is constant. Then we can do a similar thing as seen in IA Vector Calculus for the proof of uniqueness for Laplace's Equation and use the energy conservation for the difference $\psi = \phi_1 - \phi_2$ of two solutions to the Wave Equation.
\subsection{Reflection and Transmission}
We now consider singular travelling waves:
\begin{equation*}
    y(x, t) \propto \Re\left[e^{i\omega(t - x / c)}\right].
\end{equation*}
We can have a difference in density of the string for $x < 0$ and $x > 0$, in which case we need to solve on the separate domains $x > 0$ and $x < 0$, along with boundary conditions at $x = 0$ to match $y$ and $\frac{\partial y}{\partial x}$.

The other possibility is a mass at $x = 0$ (but the same density string), in which case we again solve either side of $0$ with $y$ matching at $0$, but the other matching condition requires consideration of Newton's 2nd Law (see example sheet). %TODO: Get this from the example sheet.

For both of the above setups, with incoming waves from $-\infty$, our solution looks like:
\begin{equation*}
    y(x, t) =
    \begin{cases}
        \Re\left[e^{i\omega(t - x / c_-)} + Re^{-i\omega(t - x / c_-)}\right] & x < 0 \\
        \Re\left[Re^{i\omega(t - x / c_+)}\right] & x > 0 \\
    \end{cases}
\end{equation*}
where $R$ and $T$ are the reflection and transmission coefficients.
\section{Heat Equation}
\subsection{Solution}
The heat equation is an IBVP given by:
\begin{equation*}
    \frac{\partial \phi}{\partial t} = \kappa \nabla^2 \phi
\end{equation*}
where $\kappa > 0$ is the thermal conductivity of the material.

We again solve this using separation of variables. Separating $\phi$ as $\phi(\vec{x}, t) = \psi(\vec{x}) T(t)$ gives $T = e^{-\kappa \mu t}$ for eigenvalue $\mu$, and that $\psi$ solves the Laplacian eigenvalue problem $\nabla^2 \psi = -\mu \psi$.

A square domain for $\psi$ gives sinusoidal solutions in $x$ and $y$, a cylindrical domain gives the Bessel equation in $r$ and a sinusoidal solution in $z$.
\subsection{Heat Loss and Uniqueness}
Define the heat energy:
\begin{equation*}
    Q(t) = \frac12 \int_{\Omega} \phi(\vec{x}, t)^2 dV
\end{equation*}
and by taking its derivative we see that this is a (non-strictly) decreasing function with $t$.

\begin{proposition}
    If the quantity $Q(t)$ is finite, the solution to the heat equation is unique.
    \label{propHeatUnique}
\end{proposition}
\begin{proof}
    Again consider $\psi = \phi_1 - \phi_2$ the difference of two solutions of the heat equation. Then this solution has energy $0$ at time $0$, and so must have $Q \equiv 0$ because $Q$ is non-negative and decreasing. This implies $\psi \equiv 0$.
\end{proof}
\end{document}